\subsection{The Hausdorff completion of a filtered module}

Let G be a filtered group whose filtration \((\mathbf{G}_{n})\) consists of normal subgroups of G; we have already recalled (no. 6) that the Hausdorff completion \(\hat{G}\) of the topological group G is canonically identified with the inverse limit \(\lim_{\substack{\mathbf{G}/\mathbf{G}_{n}\\ }} \overleftarrow{\mathbf{G}}\) of the discrete groups \(G/G_{n}\) , the canonical homomorphism \(i: G \to \hat{G}\) having image the Hausdorff group associated with G (everywhere dense in \(\hat{G}\) ) and kernel the closure \(\bigcap G_{n}\) of \(\{0\}\) in G. The Hausdorff completion \(\hat{G}_{n}\) of the subgroup \(G_{n}\) of G is identified with the closure of \(i(\mathbf{G}_{n})\) in \(\hat{G}\) (General Topology, Chapter II, §3, no. 9, Corollary 1 to Proposition 18) and, since \(G_{n}\) is closed in G,


\[
\mathbf {G} _ {n} = \stackrel {- 1} {i} (\hat {\mathbf {G}} _ {n}) = \stackrel {- 1} {i} (\hat {\mathbf {G}} _ {n} \cap i (\mathbf {G})).\tag{18}
\]

Moreover, the \(\hat{G}_n\) form a fundamental system of neighbourhoods of 0 in \(\hat{G}\) (General Topology, Chapter III, §3, no. 4, Proposition 7) and are therefore normal open subgroups of \(\hat{\mathbf{G}}\) (General Topology, Chapter III, § 2, no. 3, Proposition 8); the topology on \(\hat{\mathbf{G}}\) is defined by the filtration \((\hat{\mathbf{G}}_n)\) , which is always separated by definition. As \(i(\mathbf{G})\) is dense in \(\hat{\mathbf{G}}\) and \(\hat{\mathbf{G}}_n\) is open,

\[
\hat {\mathbf {G}} = i (\mathbf {G}). \hat {\mathbf {G}} _ {n}\tag{19}
\]

and similarly

\[
\hat {\mathbf {G}} _ {n - 1} = i \left(\mathbf {G} _ {n - 1}\right). \hat {\mathbf {G}} _ {n}.\tag{20}
\]

We deduce from (18) and (19) that the filtration \((\hat{G}_n)\) is exhaustive if and only if \((G_n)\) is.

The second isomorphism theorem (Algebra, Chapter I, § 6, no. 13, Theorem 6 (d)) and equations (18), (19) and (20) show that the canonical homomorphisms

\[
\mathrm{G} _ {n - 1} / \mathrm{G} _ {n} \rightarrow \hat {\mathrm{G}} _ {n - 1} / \hat {\mathrm{G}} _ {n}, \quad \mathrm{G} / \mathrm{G} _ {n} \rightarrow \hat {\mathrm{G}} / \hat {\mathrm{G}} _ {n},\tag{21}
\]

are bijective and hence so is the canonical homomorphism

\[
\operatorname{gr} (\mathbf {G}) \rightarrow \operatorname{gr} (\hat {\mathbf {G}}).\tag{22}
\]

Now let A be a filtered ring, E a filtered A-module and \((\mathbf{A}_{n})\) and \((\mathbf{E}_{n})\) the respective filtrations of A and E; we shall assume that these filtrations are exhaustive so that for the corresponding topologies A is a topological ring and E a topological A-module (no. 5, Proposition 3). Then we have defined (General Topology, Chapter III, § 6, nos. 5 and 6) \(\hat{A}\) as a topological ring and \(\hat{E}\) as a topological \(\hat{A}\) -module. If \(i: A \to \hat{A}\) is the canonical homomorphism, then \(i(\mathbf{A}_{m})i(\mathbf{A}_{n}) \subset i(\mathbf{A}_{m+n})\) , whence by the continuity of multiplication in \(\hat{A}\) ,

\[
\hat {\mathrm{A}} _ {m} \hat {\mathrm{A}} _ {n} \subset \hat {\mathrm{A}} _ {m + n}\tag{23}
\]

since \(\hat{\mathbf{A}}_n\) is the closure of \(i(\mathbf{A}_n)\) in \(\hat{\mathbf{A}}\) . It can be similarly shown that

\[
\hat {\mathrm{A}} _ {m} \hat {\mathrm{E}} _ {n} \subset \hat {\mathrm{E}} _ {m + n};\tag{24}
\]

in other words:

PROPOSITION 15. Let A be a filtered ring and E a filtered A-module, the respective filtrations \((\mathbf{A}_{n})\) , \((\mathbf{E}_{n})\) of A and E being exhaustive. Then \((\hat{\mathbf{A}}_{n})\) is a filtration compatible with the ring structure on \(\hat{A}\) and \((\hat{\mathbf{E}}_{n})\) a filtration compatible with the module structure on \(\hat{E}\) over the filtered ring \(\hat{A}\) ; moreover these filtrations are exhaustive and define respectively the topologies on \(\hat{A}\) and \(\hat{E}\) . Finally, the canonical mappings \(\mathrm{gr}(\mathbf{A}) \to \mathrm{gr}(\hat{\mathbf{A}})\) and \(\mathrm{gr}(\mathrm{E}) \to \mathrm{gr}(\hat{\mathrm{E}})\) of graded Z-modules are respectively a graded ring isomorphism and a graded gr(A)-module isomorphism.

In what follows, for every uniform space X, \(j_{x}\) will denote the canonical mapping from X to its Hausdorff completion \(\hat{X}\) and \(\mathbf{X}_{0}=j_{\mathbf{x}}(\mathbf{X})\) the uniform subspace of \(\hat{X}\) , which is the Hausdorff space associated with X. Recall that the topology on X is the inverse image under \(j_{x}\) of that on \(X_{0}\) (General Topology,

Chapter II, § 3, no. 7, Proposition 12). Recall also that, for every uniformly continuous mapping \(f\colon\mathbf{X}\to\mathbf{Y},\hat{f}\) denotes the uniformly continuous mapping from \(\hat{\mathbf{X}}\) to \(\hat{\mathbf{Y}}\) such that \(\hat{f}\circ j_{\mathbf{X}}=j_{\mathbf{Y}}\circ f\) (loc. cit., Proposition 15); if \(\mathbf{X}\) is a uniform subspace of \(\mathbf{Y}\) and \(f\) the canonical injection, \(\hat{\mathbf{X}}\) is identified with a uniform subspace of \(\hat{\mathbf{Y}}\) and \(\hat{f}\) is the canonical injection of \(\hat{\mathbf{X}}\) into \(\hat{\mathbf{Y}}\) (loc. cit., no. 9. Corollary 1 to Proposition 18).

LEMMA 2. Let \(\mathbf{X} \xrightarrow{f} \mathbf{Y} \xrightarrow{g} \mathbf{Z}\) be an exact sequence of strict morphisms of topological groups (Algebra, Chapter II, §1, no. 4, Remark). Suppose that \(\mathbf{X}, \mathbf{Y}, \mathbf{Z}\) admit Hausdorff completion groups and that the identity elements of \(\mathbf{X}, \mathbf{Y}, \mathbf{Z}\) admit countable fundamental systems of neighbourhoods. Then \(\hat{\mathbf{X}} \xrightarrow{f} \hat{\mathbf{Y}} \xrightarrow{g} \hat{\mathbf{Z}}\) is an exact sequence of strict morphisms.

Let \(N_{f}, N_{g}\) be the respective kernels of f and g; let us write

\[
f = f _ {3} \circ f _ {2} \circ f _ {1}
\]

where \(f_1\) is the canonical mapping \(\mathbf{X} \to \mathbf{X} / \mathbf{N}_f\) , \(f_2\) is an isomorphism of \(\mathbf{X} / \mathbf{N}_f\) onto \(\mathbf{N}_g\) and \(f_3\) is the canonical injection \(\mathbf{N}_g \to \mathbf{Y}\) . We already know that \(\hat{f}_2\) is an isomorphism of \((\mathbf{X} / \mathbf{N}_f)^{\wedge}\) onto \(\hat{\mathbf{N}}_g\) and we have just recalled that \(\hat{f}_3\) is an injective strict morphism of \(\hat{\mathbf{N}}_g\) to \(\hat{\mathbf{Y}}\) ; if we show that \(\hat{f}_1\) is a surjective strict morphism, it will follow that \(\hat{f}\) is a strict morphism (General Topology, Chapter III, § 2, no. 8, Remark 2). Let \(g_1\) be the canonical mapping \(\mathbf{Y} \to \mathbf{Y} / \mathbf{N}_g\) ; if we show that \(\hat{g}_1\) is a surjective strict morphism of the kernel \(\hat{\mathbf{N}}_g\) , we shall see as above that \(\hat{g}\) is a strict morphism and the sequence \(\hat{\mathbf{X}}\xrightarrow{f}\hat{\mathbf{Y}}\xrightarrow{\hat{g}}\hat{\mathbf{Z}}\) will be exact. Thus we have reduced the problem to proving that, if \(\mathbf{Y}=\mathbf{X}/\mathbf{N}\) (where \(\mathbf{N}\) is a normal subgroup of \(\mathbf{X}\) ) and \(f\colon\mathbf{X}\to\mathbf{Y}\) is the canonical mapping, \(\hat{f}\) is a surjective strict morphism with kernel \(\hat{\mathbf{N}}\) .

Let \(f_{0}: X_{0} \to Y_{0}\) be the mapping which coincides with \(\hat{f}\) on \(X_{0}\) ; as \(j_{X}\) (resp. \(j_{Y}\) ) is a surjective strict morphism of X onto \(X_{0}\) (resp. Y onto \(Y_{0}\) ), \(f_{0}\) is a surjective strict morphism (General Topology, Chapter III, § 2, no. 8, Remark 3). Now \(X_{0}\) and \(Y_{0}\) are metrizable (General Topology, Chapter IX, § 3, no. 1, Proposition 1); then it follows from General Topology, Chapter IX, § 3, no. 1, Corollary 1 to Proposition 4 and Lemma 1, that \(\hat{f}_{0} = \hat{f}\) is a surjective strict morphism and has kernel the closure \(\hat{N}_{0}'\) in \(\hat{X}\) of the kernel \(N_{0}'\) of \(f_{0}\) . Then it will be sufficient for us to prove that \(\hat{N}_{0}' = \hat{N}\) . Now \(N_{0}'\) obviously contains \(N_{0} = j_{X}(N)\) ; it will be sufficient to show that \(N_{0}'\) is contained in the closure \(\bar{N}_{0}\) of \(N_{0}\) in X. Now,

\[
\mathbf {U} = \bar {j} _ {\mathbf {X}} ^ {- 1} (\mathbf {X} _ {0} - \overline {{\mathbf {N}}} _ {0}) = \mathbf {X} - \bar {j} _ {\mathbf {X}} ^ {- 1} (\overline {{\mathbf {N}}} _ {0})
\]

is an open set in X which does not meet N; as f is a surjective strict morphism, \(V = f(U)\) is an open set in Y not containing the identity element \(e'\) of Y and hence not meeting the closure of \(e'\) ; then \(j_{Y}(V)\) does not contain the identity element of \(Y_{0}\) . But \(j_{Y}(V) = f_{0}(X_{0} - \overline{N}_{0})\) and hence \(N_{0}' \subset \overline{N}_{0}\) , which complete the proof of Lemma 2.

PROPOSITION 16. Let A be a filtered ring, \((\mathbf{A}_n)\) its filtration, E an A-module and \((\mathbf{E}_n)\) the filtration on E derived from that on A consisting of the \(\mathbf{E}_n = \mathbf{A}_n\mathbf{E}\) . Suppose that the filtration \((\mathbf{A}_n)\) is exhaustive and the module E is finitely generated. If \(i: \mathbf{E} \to \hat{\mathbf{E}}\) is the canonical mapping, then, for all \(n \in \mathbf{Z}\) ,

\[
\hat {\mathrm{E}} _ {n} = \hat {\mathrm{A}} _ {n} \hat {\mathrm{E}} = \hat {\mathrm{A}} _ {n} i (\mathrm{E}) \quad a n d \quad \hat {\mathrm{E}} = \hat {\mathrm{A}}. i (\mathrm{E}).\tag{25}
\]

In particular \(\hat{\mathbf{E}}\) is a finitely generated \(\hat{\mathbf{A}}\) -module.

The equation \(A_{n}E = E_{n}\) implies, by virtue of the continuity of the external law on the \(\hat{A}\) -module \(\hat{E}, \hat{A}_{n}\hat{E} \subset \hat{E}_{n}\) and obviously \(\hat{A}_{n}\hat{E} \supset \hat{A}_{n}i(E)\) . By hypothesis there exists a surjective homomorphism \(u: L \to E\) , where \(L = A_{s}^{I}\) , I being a finite set; let L be given the product filtration, consisting of the \(L_{n} = A_{n}^{I}\) , which define on L the product topology; then \(\hat{L} = \hat{A}_{s}^{I}\) and \(\hat{L}_{n} = \hat{A}_{n}^{I}\) (General Topology, Chapter II, § 3, no. 9, Corollary 2 to Proposition 18). Let \(j: L \to \hat{L}\) be the canonical mapping and \((e_{i})_{i \in I}\) the canonical basis of L; for an element \(\sum_{i \in I} a_{i}j(e_{i})\) (where \(a_{i} \in \hat{A}\) for all \(i \in I\) ) to belong to \(\hat{L}_{n}\) , it is necessary and sufficient that \(a_{i} \in \hat{A}_{n}\) for all i; then \(\hat{L}_{n} = \hat{A}_{n}.j(L)\) . This being so, by definition \(u(\mathbf{L}_{n}) = \mathbf{A}_{n}\mathbf{E} = \mathbf{E}_{n}\) and hence u is a strict morphism of L onto E (General Topology, Chapter III, § 2, no. 8, Proposition 24). Lemma 2 then shows that \(\hat{u}: \hat{L} \to \hat{E}\) is a surjective strict morphism. As \(\hat{L}_{n}\) is an open subgroup of \(\hat{L}\) , \(\hat{u}(\hat{\mathbf{L}}_{n})\) is an open (and therefore closed) subgroup of \(\hat{E}\) ; but \(\hat{u}(\hat{\mathbf{L}}_{n}) = \hat{\mathbf{A}}_{n}\hat{u}(j(\mathbf{L})) = \hat{\mathbf{A}}_{n}i(\mathbf{E})\) and, as \(i(\mathbf{E}_{n}) \subset \mathbf{A}_{n}i(\mathbf{E}) \subset \hat{\mathbf{A}}_{n}i(\mathbf{E})\) , finally \(\hat{\mathbf{E}}_{n} \subset \hat{\mathbf{A}}_{n}i(\mathbf{E}) \subset \hat{\mathbf{A}}_{n}\hat{\mathbf{E}} \subset \hat{\mathbf{E}}_{n}\) and therefore \(\hat{\mathbf{E}}_{n} = \hat{\mathbf{A}}_{n}\hat{\mathbf{E}} = \hat{\mathbf{A}}_{n}i(\mathbf{E})\) ; setting n = 0, we obtain the second formula of (25).

COROLLARY 1. Under the conditions of Proposition 16, if A is complete, so is E.

As the canonical mapping \(\mathbf{A} \to \hat{\mathbf{A}}\) is then surjective (no. 6, Proposition 5), \(\hat{\mathbf{E}} = i(\mathbf{E})\) by (25) and the conclusion follows by Proposition 5 of no. 6.

COROLLARY 2. Let A be a commutative ring, m a finitely generated ideal of A and \(\hat{\mathbf{A}}\) the Hausdorff completion of A with respect to the m-adic topology. Then \(\widehat{\mathfrak{m}^n} = (\hat{\mathfrak{m}})^n = \mathfrak{m}^n\) . \(\hat{\mathbf{A}}\) for every integer \(n > 0\) and the topology on \(\hat{\mathbf{A}}\) is the m-adic topology.

Let us write \(A_{n} = m^{n}\) , which is a finitely generated ideal of A. The formula \(m^{p}A_{n} = m^{n+p}\) shows that the topology induced on \(A_{n}\) by the m-adic topology coincides with the m-adic topology on the A-module \(A_{n}\) (no. 1, Example 3). By Proposition 16 applied to \(E = A_{n}\) , \(\hat{A}_{n} = \hat{A}\) . \(A_{n}\) , in other words \(\widehat{m^{n}} = m^{n}\) . \(\hat{A}\) . In particular \(\hat{m} = m\) . \(\hat{A}\) , whence

\[
(\hat {\mathfrak {m}}) ^ {n} = \mathfrak {m} ^ {n}. \hat {\mathbf {A}} = \hat {\mathbf {A}}. \mathbf {A} _ {n}
\]

(cf.~Exercise 12).

\textbf{Examples of Hausdorff completions of filtered rings}\par

\begin{enumerate}
\def\labelenumi{(\arabic{enumi})}
\tightlist
\item
  Let \(\mathbf{A}\) be a graded ring of type \(\mathbf{N}\) and let \((\mathbf{A}_n)_{n\geqslant 0}\) be its graduation; let it be given the associated filtration which is separated and exhaustive (no. 1, Example 1). The additive group \(\mathbf{A}\) is canonically identified with a subgroup of \(\mathbf{B}_{n}\) as if \(n\) is defined by the subgroup of

\(B = \prod_{n \in N} A_n\) ; if \(B\) is given the topology the product of the discrete topologies, the topology induced on \(A\) is the topology defined by the filtration on \(A\) ; also \(B\) is a complete topological group and \(A\) is dense in \(B\) (General Topology, Chapter III, § 2, no. 9, Proposition 25). The additive topological group \(B\) is then identified with the completion \(\hat{A}\) of the Hausdorff additive group \(A\) and it follows from Proposition 15 that it has a unique ring structure which makes it the completion of the topological ring \(A\) . To define multiplication in this ring,

note that, if we write \(A_{n}^{\prime} = \sum_{i > n} A_{i}\) , the closure in B of the two-sided ideal \(A_{n}^{\prime}\) is the set \(B_{n}\) of \(x = (x_{i}) \in \mathrm{B}\) such that \(x_{i} = 0\) for \(i \leqslant n\) . Then let \(x = (x_{i}), y = (y_{i})\) be two elements of B and \(z = (z_{i})\) their product. Then, for all n \textgreater{} 0, \(x \equiv x_{n}^{\prime} \pmod{B_{n}}\) , \(y \equiv y_{n}^{\prime} \pmod{B_{n}}\) , where \(x_{n}^{\prime} = (x_{i})_{0 \leqslant i \leqslant n}, y_{n}^{\prime} = (y_{i})_{0 \leqslant i \leqslant n}\) , whence \(z \equiv x_{n}^{\prime} y_{n}^{\prime} \pmod{B_{n}}\) . But \(x_{n}^{\prime}\) and \(y_{n}^{\prime}\) belong to A and it is therefore seen that, for all \(n \in N\) ,

\[
z _ {n} = \sum_ {j = 0} ^ {n} x _ {j} y _ {n - j}.\tag{26}
\]

In particular, we again obtain the Corollary to Proposition 6 of no. 6: if C is a commutative ring, the completion of the polynomial ring \(C[X_{1},\ldots,X_{r}]\) , with the filtration associated with its usual graduation (by total degree) is canonically identified with the ring of formal power series \(C[[X_{1},\ldots,X_{r}]]\) (cf.~Algebra, Chapter IV, § 5, no. 10).
\end{enumerate}

* (2) Let K be a complete commutative field with a valuation. The completion of the ring of convergent series in r variables over K is canonically identified with the ring of formal power series \(K[[X_{1},\ldots,X_{r}]]\) . *

\begin{enumerate}
\def\labelenumi{(\arabic{enumi})}
\setcounter{enumi}{2}
\tightlist
\item
  Let \(\alpha\) be a non-zero non-invertible element of a principal ideal domain; the \((\alpha)\) -adic topology on A is also called the \(\alpha\) -adic topology; it is Hausdorff, for the intersection of the ideals \((\alpha^n)\) reduces to 0 (Algebra, Chapter VII, §1, no. 3). Note that the completion of A with respect to this topology is not necessarily an integral domain (cf.~no. 13, Remark 3). The associated graded ring \(\mathrm{gr}(\mathbf{A}) = \mathrm{gr}(\hat{\mathbf{A}})\) is canonically isomorphic to \((\mathbf{A}/\mathfrak{a})[\mathbf{X}]\) (no. 3, Example 1). If \(\mathbf{A} = \mathbf{Z}\) , the completion of \(\mathbf{Z}\) with respect to the \(n\) -adic topology ( \(n > 1\) ) is denoted by \(\mathbf{Z}_n\) and its elements are called \(n\) -adic integers.
\end{enumerate}

Every element of \(\mathbf{Z} / n^k\mathbf{Z}\) admits a unique representative of the form \(\sum_{i=0}^{k-1} a_i n^i\) where \(0 \leqslant a_i \leqslant n - 1\) for all \(i\) ; moreover, its canonical image in \(\mathbf{Z} / n^{k-1}\mathbf{Z}\) is the class of \(\sum_{i=0}^{k-2} a_i n^i\) . These remarks and the fact that \(\mathbf{Z}_n\) is canonically identified with the inverse limit \(\lim_{k} \mathbf{Z}/n^k \mathbf{Z}\) show immediately that every

element of \(\mathbf{Z}_n\) can be written uniquely in the form \(\sum_{i=0}^{\infty} a_i n^i\) where \(0 \leqslant a_i < n\) and conversely that such a series is convergent in \(\mathbf{Z}_n\) .
% semantic-fragments: legacy-input-seam
\relax{}
